\section{晚明的冲突、灾害与司法证据}

战争、灾荒、疾病、赈恤、案狱和边地暴力并不是彼此独立的材料领域。军报和判牍往往由胜者、官府或审讯机关形成；灾情、医疗与慈善记录也常以请赈、捐输、工程或死亡数字的形式出现。它们能显示国家和地方组织如何命名风险与行动，却不能把受害者、病者、被征发者或被定罪者的经验全部还原。

\subsection{战事、军费与地方动员}

\eventanchor{ML-1631-005}{明·1631（崇祯四年）}
\paragraph{1631（崇祯四年）：地方结社、宗教、出版或士人文集的本年线索提供战乱中的社会观察，幸存者记忆不能概括全国。（材料核查重点：诏令或奏议的形成与发受链）} 此一账本锚点将冲突、风险或法律处置置于具体时段。它所能可靠支持的判断是来源导向的年度桥接条目：本年材料可定位相关政务线索；具体月日、发文机关、地域和执行结果仍须逐项回查，不能以制度文本或奏报替代实际效果。材料链由，因而同时保存了命令、报送、传记、地方记忆或跨政权叙事的不同立场。问题形成的条件是制度、教育或知识传播的需要推动了相关行动或文本形成；随后产生的可见影响是它提供制度和传播史的锚点，不能据此推断社会整体接受。对于医疗救助、慈善网络、普通人的伤亡与案狱处境，材料往往尤为稀薄。桥接条目只标示可回查的年度问题与材料链；实录有中央选择性，崇祯期尤其需要奏疏、地方、朝鲜及满文材料交叉。；本条特别核对诏令或奏议的形成与发受链。

\eventanchor{ML-1631-006}{明·1631（崇祯四年）}
\paragraph{1631（崇祯四年）：地方结社、宗教、出版或士人文集的本年线索提供战乱中的社会观察，幸存者记忆不能概括全国。（材料核查重点：报称信息的来源、抄传与行政分类）} 此一账本锚点将冲突、风险或法律处置置于具体时段。它所能可靠支持的判断是来源导向的年度桥接条目：本年材料可定位相关政务线索；具体月日、发文机关、地域和执行结果仍须逐项回查，不能以制度文本或奏报替代实际效果。材料链由，因而同时保存了命令、报送、传记、地方记忆或跨政权叙事的不同立场。问题形成的条件是制度、教育或知识传播的需要推动了相关行动或文本形成；随后产生的可见影响是它提供制度和传播史的锚点，不能据此推断社会整体接受。对于医疗救助、慈善网络、普通人的伤亡与案狱处境，材料往往尤为稀薄。桥接条目只标示可回查的年度问题与材料链；实录有中央选择性，崇祯期尤其需要奏疏、地方、朝鲜及满文材料交叉。；本条特别核对报称信息的来源、抄传与行政分类。

\eventanchor{ML-1631-007}{明·1631（崇祯四年）}
\paragraph{1631（崇祯四年）：天启末至崇祯朝的任免、奏疏和廷议构成本年中枢危机线索；崇祯无实录，必须以多类材料互校。（材料核查重点：同年地方材料所见的执行范围）} 此一账本锚点将冲突、风险或法律处置置于具体时段。它所能可靠支持的判断是来源导向的年度桥接条目：本年材料可定位相关政务线索；具体月日、发文机关、地域和执行结果仍须逐项回查，不能以制度文本或奏报替代实际效果。材料链由，因而同时保存了命令、报送、传记、地方记忆或跨政权叙事的不同立场。问题形成的条件是继承、官僚协调或皇权运作的压力使问题进入朝廷程序；随后产生的可见影响是它影响后续人事与政策议程；动机和派系标签不能由单一叙事裁定。对于医疗救助、慈善网络、普通人的伤亡与案狱处境，材料往往尤为稀薄。桥接条目只标示可回查的年度问题与材料链；实录有中央选择性，崇祯期尤其需要奏疏、地方、朝鲜及满文材料交叉。；本条特别核对同年地方材料所见的执行范围。

\eventanchor{ML-1632-001}{明·1632（崇祯五年）二月22日}
\paragraph{1632（崇祯五年）二月22日：吴桥兵变：山东部队兵变攻克邓州} 此一账本锚点将冲突、风险或法律处置置于具体时段。它所能可靠支持的判断是编年候选的年序可由官修与后出材料交叉；具体月日、参与者、战果或数字必须回查所列材料，不能将单一叙事当作定论。材料链由，因而同时保存了命令、报送、传记、地方记忆或跨政权叙事的不同立场。问题形成的条件是前期边防、地方武装或跨区域资源调动的压力推动了该行动；随后产生的可见影响是它改变了后续调兵、补给或地方秩序的条件；战果和伤亡须分开核对。对于医疗救助、慈善网络、普通人的伤亡与案狱处境，材料往往尤为稀薄。译写候选以编年材料定位；实录记载反映朝廷选择，崇祯期须另以奏疏、地方及敌对政权材料交叉。

\eventanchor{ML-1632-002}{明·1632（崇祯五年）}
\paragraph{1632（崇祯五年）：西班牙马尼拉与中国的贸易额每年达到200万比索} 此一账本锚点将冲突、风险或法律处置置于具体时段。它所能可靠支持的判断是编年候选的年序可由官修与后出材料交叉；具体月日、参与者、战果或数字必须回查所列材料，不能将单一叙事当作定论。材料链由，因而同时保存了命令、报送、传记、地方记忆或跨政权叙事的不同立场。问题形成的条件是朝贡名义、边境安全与港口贸易的利益使交涉持续发生；随后产生的可见影响是它为后续册封、互市或海防安排留下文书与跨政权比较线索。对于医疗救助、慈善网络、普通人的伤亡与案狱处境，材料往往尤为稀薄。译写候选以编年材料定位；实录记载反映朝廷选择，崇祯期须另以奏疏、地方及敌对政权材料交叉。

\eventanchor{ML-1632-003}{明·1632（崇祯五年）}
\paragraph{1632（崇祯五年）：明朝的防御规划者建造了一些星堡，但在中国并没有流行起来} 此一账本锚点将冲突、风险或法律处置置于具体时段。它所能可靠支持的判断是编年候选的年序可由官修与后出材料交叉；具体月日、参与者、战果或数字必须回查所列材料，不能将单一叙事当作定论。材料链由，因而同时保存了命令、报送、传记、地方记忆或跨政权叙事的不同立场。问题形成的条件是继承、官僚协调或皇权运作的压力使问题进入朝廷程序；随后产生的可见影响是它影响后续人事与政策议程；动机和派系标签不能由单一叙事裁定。对于医疗救助、慈善网络、普通人的伤亡与案狱处境，材料往往尤为稀薄。译写候选以编年材料定位；实录记载反映朝廷选择，崇祯期须另以奏疏、地方及敌对政权材料交叉。

\subsection{灾害、医疗与赈恤}

\eventanchor{ML-1632-004}{明·1632（崇祯五年）}
\paragraph{1632（崇祯五年）：辽饷、剿饷、练饷及地方征解的本年文书显示财政负担，定额、征收与实际流向不得混为一谈。（材料核查重点：诏令或奏议的形成与发受链）} 此一账本锚点将冲突、风险或法律处置置于具体时段。它所能可靠支持的判断是来源导向的年度桥接条目：本年材料可定位相关政务线索；具体月日、发文机关、地域和执行结果仍须逐项回查，不能以制度文本或奏报替代实际效果。材料链由，因而同时保存了命令、报送、传记、地方记忆或跨政权叙事的不同立场。问题形成的条件是财政征解、交通工程或货币支付的压力使相关措施进入政务；随后产生的可见影响是其法定安排不等于地方执行，后续须以地域材料追踪负担与成效。对于医疗救助、慈善网络、普通人的伤亡与案狱处境，材料往往尤为稀薄。桥接条目只标示可回查的年度问题与材料链；实录有中央选择性，崇祯期尤其需要奏疏、地方、朝鲜及满文材料交叉。；本条特别核对诏令或奏议的形成与发受链。

\eventanchor{ML-1632-005}{明·1632（崇祯五年）}
\paragraph{1632（崇祯五年）：旱灾、蝗灾、疫病、河患与赈济的本年材料连接环境和社会压力，死亡与流离数字须谨慎处理。（材料核查重点：诏令或奏议的形成与发受链）} 此一账本锚点将冲突、风险或法律处置置于具体时段。它所能可靠支持的判断是来源导向的年度桥接条目：本年材料可定位相关政务线索；具体月日、发文机关、地域和执行结果仍须逐项回查，不能以制度文本或奏报替代实际效果。材料链由，因而同时保存了命令、报送、传记、地方记忆或跨政权叙事的不同立场。问题形成的条件是自然风险与地方报告促成朝廷或地方的应对记录；随后产生的可见影响是赈济、迁徙和损失的实际范围仍须以受灾地区材料分别核验。对于医疗救助、慈善网络、普通人的伤亡与案狱处境，材料往往尤为稀薄。桥接条目只标示可回查的年度问题与材料链；实录有中央选择性，崇祯期尤其需要奏疏、地方、朝鲜及满文材料交叉。；本条特别核对诏令或奏议的形成与发受链。

\eventanchor{ML-1632-006}{明·1632（崇祯五年）}
\paragraph{1632（崇祯五年）：旱灾、蝗灾、疫病、河患与赈济的本年材料连接环境和社会压力，死亡与流离数字须谨慎处理。（材料核查重点：报称信息的来源、抄传与行政分类）} 此一账本锚点将冲突、风险或法律处置置于具体时段。它所能可靠支持的判断是来源导向的年度桥接条目：本年材料可定位相关政务线索；具体月日、发文机关、地域和执行结果仍须逐项回查，不能以制度文本或奏报替代实际效果。材料链由，因而同时保存了命令、报送、传记、地方记忆或跨政权叙事的不同立场。问题形成的条件是自然风险与地方报告促成朝廷或地方的应对记录；随后产生的可见影响是赈济、迁徙和损失的实际范围仍须以受灾地区材料分别核验。对于医疗救助、慈善网络、普通人的伤亡与案狱处境，材料往往尤为稀薄。桥接条目只标示可回查的年度问题与材料链；实录有中央选择性，崇祯期尤其需要奏疏、地方、朝鲜及满文材料交叉。；本条特别核对报称信息的来源、抄传与行政分类。

\eventanchor{ML-1632-007}{明·1632（崇祯五年）}
\paragraph{1632（崇祯五年）：辽饷、剿饷、练饷及地方征解的本年文书显示财政负担，定额、征收与实际流向不得混为一谈。（材料核查重点：报称信息的来源、抄传与行政分类）} 此一账本锚点将冲突、风险或法律处置置于具体时段。它所能可靠支持的判断是来源导向的年度桥接条目：本年材料可定位相关政务线索；具体月日、发文机关、地域和执行结果仍须逐项回查，不能以制度文本或奏报替代实际效果。材料链由，因而同时保存了命令、报送、传记、地方记忆或跨政权叙事的不同立场。问题形成的条件是财政征解、交通工程或货币支付的压力使相关措施进入政务；随后产生的可见影响是其法定安排不等于地方执行，后续须以地域材料追踪负担与成效。对于医疗救助、慈善网络、普通人的伤亡与案狱处境，材料往往尤为稀薄。桥接条目只标示可回查的年度问题与材料链；实录有中央选择性，崇祯期尤其需要奏疏、地方、朝鲜及满文材料交叉。；本条特别核对报称信息的来源、抄传与行政分类。

\eventanchor{ML-1633-001}{明·1633（崇祯六年）四月}
\paragraph{1633（崇祯六年）四月：吴桥兵变：山东叛军投奔后金} 此一账本锚点将冲突、风险或法律处置置于具体时段。它所能可靠支持的判断是编年候选的年序可由官修与后出材料交叉；具体月日、参与者、战果或数字必须回查所列材料，不能将单一叙事当作定论。材料链由，因而同时保存了命令、报送、传记、地方记忆或跨政权叙事的不同立场。问题形成的条件是前期边防、地方武装或跨区域资源调动的压力推动了该行动；随后产生的可见影响是它改变了后续调兵、补给或地方秩序的条件；战果和伤亡须分开核对。对于医疗救助、慈善网络、普通人的伤亡与案狱处境，材料往往尤为稀薄。译写候选以编年材料定位；实录记载反映朝廷选择，崇祯期须另以奏疏、地方及敌对政权材料交叉。

\eventanchor{ML-1633-002}{明·1633（崇祯六年）七月7日}
\paragraph{1633（崇祯六年）七月7日：辽洛湾之战：明朝造船厂在郑芝龙的监督下开始建造能够容纳大炮的多层舷侧帆船；他们被荷兰人的突然袭击炸毁} 此一账本锚点将冲突、风险或法律处置置于具体时段。它所能可靠支持的判断是编年候选的年序可由官修与后出材料交叉；具体月日、参与者、战果或数字必须回查所列材料，不能将单一叙事当作定论。材料链由，因而同时保存了命令、报送、传记、地方记忆或跨政权叙事的不同立场。问题形成的条件是朝贡名义、边境安全与港口贸易的利益使交涉持续发生；随后产生的可见影响是它为后续册封、互市或海防安排留下文书与跨政权比较线索。对于医疗救助、慈善网络、普通人的伤亡与案狱处境，材料往往尤为稀薄。译写候选以编年材料定位；实录记载反映朝廷选择，崇祯期须另以奏疏、地方及敌对政权材料交叉。

\subsection{审案、处分与法律语言}

\eventanchor{ML-1633-003}{明·1633（崇祯六年）夏}
\paragraph{1633（崇祯六年）夏：围攻旅顺：后金占领旅顺} 此一账本锚点将冲突、风险或法律处置置于具体时段。它所能可靠支持的判断是编年候选的年序可由官修与后出材料交叉；具体月日、参与者、战果或数字必须回查所列材料，不能将单一叙事当作定论。材料链由，因而同时保存了命令、报送、传记、地方记忆或跨政权叙事的不同立场。问题形成的条件是前期边防、地方武装或跨区域资源调动的压力推动了该行动；随后产生的可见影响是它改变了后续调兵、补给或地方秩序的条件；战果和伤亡须分开核对。对于医疗救助、慈善网络、普通人的伤亡与案狱处境，材料往往尤为稀薄。译写候选以编年材料定位；实录记载反映朝廷选择，崇祯期须另以奏疏、地方及敌对政权材料交叉。

\eventanchor{ML-1633-004}{明·1633（崇祯六年）十月22日}
\paragraph{1633（崇祯六年）十月22日：辽罗湾海战：明军在金门岛附近击败荷兰海盗船队} 此一账本锚点将冲突、风险或法律处置置于具体时段。它所能可靠支持的判断是编年候选的年序可由官修与后出材料交叉；具体月日、参与者、战果或数字必须回查所列材料，不能将单一叙事当作定论。材料链由，因而同时保存了命令、报送、传记、地方记忆或跨政权叙事的不同立场。问题形成的条件是朝贡名义、边境安全与港口贸易的利益使交涉持续发生；随后产生的可见影响是它为后续册封、互市或海防安排留下文书与跨政权比较线索。对于医疗救助、慈善网络、普通人的伤亡与案狱处境，材料往往尤为稀薄。译写候选以编年材料定位；实录记载反映朝廷选择，崇祯期须另以奏疏、地方及敌对政权材料交叉。

\eventanchor{ML-1633-005}{明·1633（崇祯六年）十二月27日}
\paragraph{1633（崇祯六年）十二月27日：叛军占领绵芝} 此一账本锚点将冲突、风险或法律处置置于具体时段。它所能可靠支持的判断是编年候选的年序可由官修与后出材料交叉；具体月日、参与者、战果或数字必须回查所列材料，不能将单一叙事当作定论。材料链由，因而同时保存了命令、报送、传记、地方记忆或跨政权叙事的不同立场。问题形成的条件是继承、官僚协调或皇权运作的压力使问题进入朝廷程序；随后产生的可见影响是它影响后续人事与政策议程；动机和派系标签不能由单一叙事裁定。对于医疗救助、慈善网络、普通人的伤亡与案狱处境，材料往往尤为稀薄。译写候选以编年材料定位；实录记载反映朝廷选择，崇祯期须另以奏疏、地方及敌对政权材料交叉。

\eventanchor{ML-1633-006}{明·1633（崇祯六年）}
\paragraph{1633（崇祯六年）：天启末至崇祯朝的任免、奏疏和廷议构成本年中枢危机线索；崇祯无实录，必须以多类材料互校} 此一账本锚点将冲突、风险或法律处置置于具体时段。它所能可靠支持的判断是来源导向的年度桥接条目：本年材料可定位相关政务线索；具体月日、发文机关、地域和执行结果仍须逐项回查，不能以制度文本或奏报替代实际效果。材料链由，因而同时保存了命令、报送、传记、地方记忆或跨政权叙事的不同立场。问题形成的条件是继承、官僚协调或皇权运作的压力使问题进入朝廷程序；随后产生的可见影响是它影响后续人事与政策议程；动机和派系标签不能由单一叙事裁定。对于医疗救助、慈善网络、普通人的伤亡与案狱处境，材料往往尤为稀薄。桥接条目只标示可回查的年度问题与材料链；实录有中央选择性，崇祯期尤其需要奏疏、地方、朝鲜及满文材料交叉。

\eventanchor{ML-1633-007}{明·1633（崇祯六年）}
\paragraph{1633（崇祯六年）：辽东防务、后金（清）军事行动和流动作战集团的本年记录标记多战场互动，军数和胜败须保留分歧} 此一账本锚点将冲突、风险或法律处置置于具体时段。它所能可靠支持的判断是来源导向的年度桥接条目：本年材料可定位相关政务线索；具体月日、发文机关、地域和执行结果仍须逐项回查，不能以制度文本或奏报替代实际效果。材料链由，因而同时保存了命令、报送、传记、地方记忆或跨政权叙事的不同立场。问题形成的条件是前期边防、地方武装或跨区域资源调动的压力推动了该行动；随后产生的可见影响是它改变了后续调兵、补给或地方秩序的条件；战果和伤亡须分开核对。对于医疗救助、慈善网络、普通人的伤亡与案狱处境，材料往往尤为稀薄。桥接条目只标示可回查的年度问题与材料链；实录有中央选择性，崇祯期尤其需要奏疏、地方、朝鲜及满文材料交叉。

\eventanchor{ML-1634-001}{明·1634（崇祯七年）九月14日}
\paragraph{1634（崇祯七年）九月14日：桐城叛乱爆发} 此一账本锚点将冲突、风险或法律处置置于具体时段。它所能可靠支持的判断是编年候选的年序可由官修与后出材料交叉；具体月日、参与者、战果或数字必须回查所列材料，不能将单一叙事当作定论。材料链由，因而同时保存了命令、报送、传记、地方记忆或跨政权叙事的不同立场。问题形成的条件是前期边防、地方武装或跨区域资源调动的压力推动了该行动；随后产生的可见影响是它改变了后续调兵、补给或地方秩序的条件；战果和伤亡须分开核对。对于医疗救助、慈善网络、普通人的伤亡与案狱处境，材料往往尤为稀薄。译写候选以编年材料定位；实录记载反映朝廷选择，崇祯期须另以奏疏、地方及敌对政权材料交叉。

\subsection{边地暴力与行政后果}

\eventanchor{ML-1634-002}{明·1634（崇祯七年）}
\paragraph{1634（崇祯七年）：地方结社、宗教、出版或士人文集的本年线索提供战乱中的社会观察，幸存者记忆不能概括全国} 此一账本锚点将冲突、风险或法律处置置于具体时段。它所能可靠支持的判断是来源导向的年度桥接条目：本年材料可定位相关政务线索；具体月日、发文机关、地域和执行结果仍须逐项回查，不能以制度文本或奏报替代实际效果。材料链由，因而同时保存了命令、报送、传记、地方记忆或跨政权叙事的不同立场。问题形成的条件是制度、教育或知识传播的需要推动了相关行动或文本形成；随后产生的可见影响是它提供制度和传播史的锚点，不能据此推断社会整体接受。对于医疗救助、慈善网络、普通人的伤亡与案狱处境，材料往往尤为稀薄。桥接条目只标示可回查的年度问题与材料链；实录有中央选择性，崇祯期尤其需要奏疏、地方、朝鲜及满文材料交叉。

\eventanchor{ML-1634-003}{明·1634（崇祯七年）}
\paragraph{1634（崇祯七年）：辽饷、剿饷、练饷及地方征解的本年文书显示财政负担，定额、征收与实际流向不得混为一谈。（材料核查重点：诏令或奏议的形成与发受链）} 此一账本锚点将冲突、风险或法律处置置于具体时段。它所能可靠支持的判断是来源导向的年度桥接条目：本年材料可定位相关政务线索；具体月日、发文机关、地域和执行结果仍须逐项回查，不能以制度文本或奏报替代实际效果。材料链由，因而同时保存了命令、报送、传记、地方记忆或跨政权叙事的不同立场。问题形成的条件是财政征解、交通工程或货币支付的压力使相关措施进入政务；随后产生的可见影响是其法定安排不等于地方执行，后续须以地域材料追踪负担与成效。对于医疗救助、慈善网络、普通人的伤亡与案狱处境，材料往往尤为稀薄。桥接条目只标示可回查的年度问题与材料链；实录有中央选择性，崇祯期尤其需要奏疏、地方、朝鲜及满文材料交叉。；本条特别核对诏令或奏议的形成与发受链。

\eventanchor{ML-1634-004}{明·1634（崇祯七年）}
\paragraph{1634（崇祯七年）：辽东防务、后金（清）军事行动和流动作战集团的本年记录标记多战场互动，军数和胜败须保留分歧} 此一账本锚点将冲突、风险或法律处置置于具体时段。它所能可靠支持的判断是来源导向的年度桥接条目：本年材料可定位相关政务线索；具体月日、发文机关、地域和执行结果仍须逐项回查，不能以制度文本或奏报替代实际效果。材料链由，因而同时保存了命令、报送、传记、地方记忆或跨政权叙事的不同立场。问题形成的条件是前期边防、地方武装或跨区域资源调动的压力推动了该行动；随后产生的可见影响是它改变了后续调兵、补给或地方秩序的条件；战果和伤亡须分开核对。对于医疗救助、慈善网络、普通人的伤亡与案狱处境，材料往往尤为稀薄。桥接条目只标示可回查的年度问题与材料链；实录有中央选择性，崇祯期尤其需要奏疏、地方、朝鲜及满文材料交叉。

\eventanchor{ML-1634-005}{明·1634（崇祯七年）}
\paragraph{1634（崇祯七年）：天启末至崇祯朝的任免、奏疏和廷议构成本年中枢危机线索；崇祯无实录，必须以多类材料互校} 此一账本锚点将冲突、风险或法律处置置于具体时段。它所能可靠支持的判断是来源导向的年度桥接条目：本年材料可定位相关政务线索；具体月日、发文机关、地域和执行结果仍须逐项回查，不能以制度文本或奏报替代实际效果。材料链由，因而同时保存了命令、报送、传记、地方记忆或跨政权叙事的不同立场。问题形成的条件是继承、官僚协调或皇权运作的压力使问题进入朝廷程序；随后产生的可见影响是它影响后续人事与政策议程；动机和派系标签不能由单一叙事裁定。对于医疗救助、慈善网络、普通人的伤亡与案狱处境，材料往往尤为稀薄。桥接条目只标示可回查的年度问题与材料链；实录有中央选择性，崇祯期尤其需要奏疏、地方、朝鲜及满文材料交叉。

\eventanchor{ML-1634-006}{明·1634（崇祯七年）}
\paragraph{1634（崇祯七年）：辽饷、剿饷、练饷及地方征解的本年文书显示财政负担，定额、征收与实际流向不得混为一谈。（材料核查重点：报称信息的来源、抄传与行政分类）} 此一账本锚点将冲突、风险或法律处置置于具体时段。它所能可靠支持的判断是来源导向的年度桥接条目：本年材料可定位相关政务线索；具体月日、发文机关、地域和执行结果仍须逐项回查，不能以制度文本或奏报替代实际效果。材料链由，因而同时保存了命令、报送、传记、地方记忆或跨政权叙事的不同立场。问题形成的条件是财政征解、交通工程或货币支付的压力使相关措施进入政务；随后产生的可见影响是其法定安排不等于地方执行，后续须以地域材料追踪负担与成效。对于医疗救助、慈善网络、普通人的伤亡与案狱处境，材料往往尤为稀薄。桥接条目只标示可回查的年度问题与材料链；实录有中央选择性，崇祯期尤其需要奏疏、地方、朝鲜及满文材料交叉。；本条特别核对报称信息的来源、抄传与行政分类。

\eventanchor{ML-1634-007}{明·1634（崇祯七年）}
\paragraph{1634（崇祯七年）：朝鲜、辽东和海上关系的本年材料提供外部观察位置，明、后金（清）与地方势力的称谓需具体化} 此一账本锚点将冲突、风险或法律处置置于具体时段。它所能可靠支持的判断是来源导向的年度桥接条目：本年材料可定位相关政务线索；具体月日、发文机关、地域和执行结果仍须逐项回查，不能以制度文本或奏报替代实际效果。材料链由，因而同时保存了命令、报送、传记、地方记忆或跨政权叙事的不同立场。问题形成的条件是朝贡名义、边境安全与港口贸易的利益使交涉持续发生；随后产生的可见影响是它为后续册封、互市或海防安排留下文书与跨政权比较线索。对于医疗救助、慈善网络、普通人的伤亡与案狱处境，材料往往尤为稀薄。桥接条目只标示可回查的年度问题与材料链；实录有中央选择性，崇祯期尤其需要奏疏、地方、朝鲜及满文材料交叉。
